\documentclass{article}
\usepackage{amsmath,amssymb,amsthm}
\usepackage{algorithm}
\usepackage[noend]{algpseudocode}
\usepackage{bm}
\usepackage{hyperref}
\usepackage{enumitem}
\newtheorem{theorem}{Theorem}
\newtheorem{lemma}{Lemma}
\newtheorem{assumption}{Assumption}
\newcommand{\E}{\mathbb{E}}
\newcommand{\R}{\mathbb{R}}

\begin{document}

\section*{Quantized Stochastic Gradient Descent (QSGD)}

\section*{Mathematical Formulation}
Let $f:\R^{d}\rightarrow \R$ be a convex and $L$‑smooth objective.  
At iteration $t$ we have a current iterate $w_{t}\in\R^{d}$ and a (possibly stochastic) gradient $g_{t}:=\nabla f(w_{t})$ (or an unbiased stochastic estimate thereof).  
A (random) quantization operator $Q:\R^{d}\rightarrow\R^{d}$ satisfies  

\[
\E[\,Q(g)\mid g\,]=g,\qquad 
\E\big[\|Q(g)-g\|^{2}\mid g\big]\leq \beta\,\|g\|^{2},
\tag{1}
\]

for some constant $\beta\ge 0$ that quantifies the distortion introduced by quantization.  

The **QSGD update** with stepsize $\eta_{t}>0$ is  

\[
\boxed{\,w_{t+1}=w_{t}-\eta_{t}\,Q\!\big(g_{t}\big)\,}. \tag{2}
\]

The algorithm communicates only the quantized vector $Q(g_{t})$ to the parameter server (or neighboring workers), thereby reducing communication bandwidth.

\section*{Pseudocode}
\begin{algorithm}[H]
\caption{Quantized Stochastic Gradient Descent}
\begin{algorithmic}[1]
\Require Initial point $w_{0}\in\R^{d}$, stepsize schedule $\{\eta_{t}\}_{t\ge 0}$, quantizer $Q(\cdot)$ satisfying (1)
\For{$t=0,1,\dots,T-1$}
    \State Compute (or sample) a stochastic gradient $g_{t}$ at $w_{t}$ 
    \State Quantize: $\tilde g_{t}=Q(g_{t})$   \Comment{communicate $\tilde g_{t}$}
    \State Update: $w_{t+1}=w_{t}-\eta_{t}\tilde g_{t}$
\EndFor
\State \Return $w_{T}$ (or the average $\bar w_{T}:=\frac{1}{T}\sum_{t=0}^{T-1}w_{t}$)
\end{algorithmic}
\end{algorithm}

\section*{Dry Run Example}
Consider the one‑dimensional quadratic  

\[
f(w)=(w-3)^{2},\qquad \nabla f(w)=2(w-3).
\]

Take a deterministic “quantizer’’ that simply adds zero‑mean noise $ \xi_{t}$ with $\E[\xi_{t}]=0$ and $\E[\xi_{t}^{2}]=\beta\,\|\nabla f(w_{t})\|^{2}$ (so (1) holds).  
Set $w_{0}=0$, constant stepsize $\eta=0.1$, and run three iterations.

\begin{center}
\begin{tabular}{c|c|c|c|c}
$t$ & $w_{t}$ & $g_{t}=2(w_{t}-3)$ & $\tilde g_{t}=g_{t}+\xi_{t}$ & $w_{t+1}=w_{t}-\eta\tilde g_{t}$\\ \hline
0 & $0$ & $-6$ & $-6+\xi_{0}$ & $0-0.1(-6+\xi_{0})=0.6-0.1\xi_{0}$\\
1 & $0.6-0.1\xi_{0}$ & $2(0.6-0.1\xi_{0}-3)=-4.8-0.2\xi_{0}$ & $-4.8-0.2\xi_{0}+\xi_{1}$ & $0.6-0.1\xi_{0}-0.1(-4.8-0.2\xi_{0}+\xi_{1})$\\
2 & $1.08-0.08\xi_{0}-0.1\xi_{1}$ & $2(1.08-0.08\xi_{0}-0.1\xi_{1}-3)=-3.84-0.16\xi_{0}-0.2\xi_{1}$ & $-3.84-0.16\xi_{0}-0.2\xi_{1}+\xi_{2}$ & $1.08-0.08\xi_{0}-0.1\xi_{1}-0.1(-3.84-0.16\xi_{0}-0.2\xi_{1}+\xi_{2})$
\end{tabular}
\end{center}

Taking expectations (all $\xi_{t}$ have zero mean) yields the deterministic trajectory  

\[
\begin{aligned}
\E[w_{1}] &=0.6,\\
\E[w_{2}] &=0.6+0.48=1.08,\\
\E[w_{3}] &=1.08+0.384=1.464,
\end{aligned}
\]

and the corresponding objective values $\E[f(w_{t})]$ decrease accordingly, illustrating the effect of quantized updates.

\newpage
\section*{Assumptions}
\begin{assumption}[Convexity and Smoothness]\label{ass:convex}
The objective $f$ is convex and $L$‑smooth, i.e. for all $x,y\in\R^{d}$,
\[
f(y)\le f(x)+\langle\nabla f(x),y-x\rangle+\frac{L}{2}\|y-x\|^{2}.
\]
\end{assumption}

\begin{assumption}[Unbiased Stochastic Gradient]\label{ass:unbiased}
For each $t$, the (possibly stochastic) gradient $g_{t}$ satisfies
\[
\E[g_{t}\mid w_{t}]=\nabla f(w_{t}),\qquad
\E[\|g_{t}-\nabla f(w_{t})\|^{2}\mid w_{t}]\le\sigma^{2}.
\]
\end{assumption}

\begin{assumption}[Quantizer]\label{ass:quant}
The random mapping $Q:\R^{d}\to\R^{d}$ is independent of $w_{t}$ given $g_{t}$ and fulfills (1) with a constant $\beta\ge0$.
\end{assumption}

\begin{assumption}[Stepsize]\label{ass:stepsize}
The stepsizes are non‑increasing and satisfy
\[
\eta_{t}= \frac{c}{\sqrt{t+1}}\quad\text{for some }c>0,\qquad 
\sum_{t=0}^{\infty}\eta_{t}= \infty,\;\; \sum_{t=0}^{\infty}\eta_{t}^{2}<\infty .
\]
\end{assumption}

\section*{Main Convergence Theorem}
\begin{theorem}[Expected Suboptimality of QSGD]\label{thm:main}
Under Assumptions \ref{ass:convex}–\ref{ass:stepsize}, let $w^{\star}\in\arg\min f$ and define the average iterate $\bar w_{T}:=\frac{1}{T}\sum_{t=0}^{T-1}w_{t}$. Then for any $T\ge1$,
\[
\E\!\big[f(\bar w_{T})-f(w^{\star})\big]\;\le\;
\frac{\|w_{0}-w^{\star}\|^{2}}{2c\sqrt{T}}
\;+\;
\frac{c(L+\beta L+\beta\sigma^{2})}{\sqrt{T}} .
\]
Consequently, $\E[f(\bar w_{T})-f(w^{\star})]=\mathcal O(T^{-1/2})$.
\end{theorem}

\section*{Theoretical Properties}
\begin{itemize}[leftmargin=*]
\item \textbf{Stability.} The unbiasedness (1) guarantees that the quantization error does not introduce systematic drift; the variance bound $\beta$ only inflates the effective noise level.
\item \textbf{Hyper‑parameter Sensitivity.} The constant $c$ in the stepsize schedule trades off the $1/\sqrt{T}$ term against the variance term. Larger $c$ accelerates early progress but amplifies the steady‑state error proportional to $\beta$.
\item \textbf{Comparison to Plain SGD.} Setting $\beta=0$ recovers the classical SGD bound (see Example 1). The extra factor $(1+\beta)$ in the rate quantifies the communication‑induced penalty.
\item \textbf{Special Cases.} If $Q$ is a deterministic rounding with $\beta=0$, QSGD reduces to exact SGD; if $Q$ is a stochastic sign compressor with $\beta=1$, the bound becomes the known $\mathcal O(1/\sqrt{T})$ rate for sign‑SGD.
\end{itemize}

\section*{Discussion}
The theorem shows that, despite transmitting only a compressed version of the gradient, the algorithm preserves the optimal $\mathcal O(T^{-1/2})$ convergence order for convex smooth objectives. The price paid is a multiplicative constant depending linearly on the quantization variance parameter $\beta$. This matches the intuition that unbiased compression does not affect asymptotic rates, only the prefactor.

\newpage
\section*{Preliminaries (Definitions and Lemmas)}
\begin{lemma}[Smoothness Inequality]\label{lem:smooth}
If $f$ is $L$‑smooth, then for any $x,y\in\R^{d}$,
\[
f(y)\le f(x)+\langle\nabla f(x),y-x\rangle+\frac{L}{2}\|y-x\|^{2}.
\]
\end{lemma}
\begin{proof}
Standard consequence of $L$‑smoothness; see e.g. \cite{nesterov2004introductory}.
\end{proof}

\begin{lemma}[One‑step Descent Expectation]\label{lem:onestep}
Let $w_{t+1}=w_{t}-\eta_{t}Q(g_{t})$ with $g_{t}$ satisfying Assumption~\ref{ass:unbiased} and $Q$ satisfying Assumption~\ref{ass:quant}. Then
\[
\E\!\big[\|w_{t+1}-w^{\star}\|^{2}\mid \mathcal{F}_{t}\big]
=
\|w_{t}-w^{\star}\|^{2}
-2\eta_{t}\langle\nabla f(w_{t}),w_{t}-w^{\star}\rangle
+\eta_{t}^{2}\big(\|\nabla f(w_{t})\|^{2}+\beta\|\nabla f(w_{t})\|^{2}+\sigma^{2}\big),
\]
where $\mathcal{F}_{t}$ denotes the sigma‑algebra generated by $\{w_{0},\dots,w_{t}\}$.
\end{lemma}
\begin{proof}
Condition on $\mathcal{F}_{t}$. Expand the squared norm:
\[
\|w_{t+1}-w^{\star}\|^{2}
=\|w_{t}-w^{\star}\|^{2}
-2\eta_{t}\langle Q(g_{t}),w_{t}-w^{\star}\rangle
+\eta_{t}^{2}\|Q(g_{t})\|^{2}.
\tag{3}
\]
Take expectation and use unbiasedness of $g_{t}$ and $Q$:
\[
\E[Q(g_{t})\mid\mathcal{F}_{t}]
=\E[\E[Q(g_{t})\mid g_{t}]\mid\mathcal{F}_{t}]
=\E[g_{t}\mid\mathcal{F}_{t}]
=\nabla f(w_{t}).
\tag{4}
\]
Hence
\[
\E\!\big[\langle Q(g_{t}),w_{t}-w^{\star}\rangle\mid\mathcal{F}_{t}\big]
=\langle\nabla f(w_{t}),w_{t}-w^{\star}\rangle.
\tag{5}
\]
For the quadratic term,
\[
\begin{aligned}
\E\!\big[\|Q(g_{t})\|^{2}\mid\mathcal{F}_{t}\big]
&=\E\!\big[\|Q(g_{t})-g_{t}+g_{t}\|^{2}\mid\mathcal{F}_{t}\big] \\
&=\E\!\big[\|g_{t}\|^{2}\mid\mathcal{F}_{t}\big]
+\E\!\big[\|Q(g_{t})-g_{t}\|^{2}\mid\mathcal{F}_{t}\big] \\
&\le \|\nabla f(w_{t})\|^{2}+\sigma^{2}
+\beta\|\nabla f(w_{t})\|^{2} \qquad\text{(by Assumptions~\ref{ass:unbiased},~\ref{ass:quant})}.
\end{aligned}
\tag{6}
\]
Plugging (5) and (6) into (3) yields the claimed identity. ∎
\end{proof}

\section*{Main Proof (Step‑by‑Step Derivation)}
We prove Theorem~\ref{thm:main}. Throughout, $\mathcal{F}_{t}$ denotes the filtration up to iteration $t$.

\bigskip
\noindent\textbf{[Step 1]} (Smoothness upper bound). Using Lemma~\ref{lem:smooth} with $x=w_{t}$ and $y=w_{t+1}$,
\[
f(w_{t+1})\le f(w_{t})+\langle\nabla f(w_{t}),w_{t+1}-w_{t}\rangle+\frac{L}{2}\|w_{t+1}-w_{t}\|^{2}.
\tag{7}
\]

\noindent\textbf{[Step 2]} (Insert the update (2)). Since $w_{t+1}-w_{t}=-\eta_{t}Q(g_{t})$,
\[
\begin{aligned}
f(w_{t+1}) &\le f(w_{t})-\eta_{t}\langle\nabla f(w_{t}),Q(g_{t})\rangle
+\frac{L\eta_{t}^{2}}{2}\|Q(g_{t})\|^{2}.
\end{aligned}
\tag{8}
\]

\noindent\textbf{[Step 3]} (Take conditional expectation). Apply $\E[\cdot\mid\mathcal{F}_{t}]$ and use (4)–(6):
\[
\begin{aligned}
\E\!\big[f(w_{t+1})\mid\mathcal{F}_{t}\big]
&\le f(w_{t})-\eta_{t}\langle\nabla f(w_{t}),\nabla f(w_{t})\rangle
+\frac{L\eta_{t}^{2}}{2}\Big(\|\nabla f(w_{t})\|^{2}+\beta\|\nabla f(w_{t})\|^{2}+\sigma^{2}\Big)\\
&= f(w_{t})-\eta_{t}\|\nabla f(w_{t})\|^{2}
+\frac{L\eta_{t}^{2}}{2}\big((1+\beta)\|\nabla f(w_{t})\|^{2}+\sigma^{2}\big).
\end{aligned}
\tag{9}
\]

\noindent\textbf{[Step 4]} (Convexity lower bound). Convexity gives
\[
f(w_{t})-f(w^{\star})\le\langle\nabla f(w_{t}),w_{t}-w^{\star}\rangle.
\tag{10}
\]

\noindent\textbf{[Step 5]} (Relate $\|\nabla f(w_{t})\|^{2}$ to distance). By the standard inequality derived from (10) and Lemma~\ref{lem:onestep},
\[
\begin{aligned}
2\eta_{t}\big(f(w_{t})-f(w^{\star})\big)
&\le 2\eta_{t}\langle\nabla f(w_{t}),w_{t}-w^{\star}\rangle \\
&= \|w_{t}-w^{\star}\|^{2}-\E\!\big[\|w_{t+1}-w^{\star}\|^{2}\mid\mathcal{F}_{t}\big]
+\eta_{t}^{2}\big((1+\beta)\|\nabla f(w_{t})\|^{2}+\sigma^{2}\big) .
\end{aligned}
\tag{11}
\]

\noindent\textbf{[Step 6]} (Combine (9) and (11)). Rearranging (11) gives
\[
\eta_{t}\|\nabla f(w_{t})\|^{2}
\le \frac{1}{2}\Big(\|w_{t}-w^{\star}\|^{2}
-\E\!\big[\|w_{t+1}-w^{\star}\|^{2}\mid\mathcal{F}_{t}\big]\Big)
+\frac{\eta_{t}^{2}}{2}\big((1+\beta)\|\nabla f(w_{t})\|^{2}+\sigma^{2}\big).
\tag{12}
\]

\noindent\textbf{[Step 7]} (Plug (12) into (9)}. Substitute the bound on $\eta_{t}\|\nabla f(w_{t})\|^{2}$:
\[
\begin{aligned}
\E\!\big[f(w_{t+1})\mid\mathcal{F}_{t}\big]
&\le f(w_{t})-\Bigg\{\frac{1}{2}\Big(\|w_{t}-w^{\star}\|^{2}
-\E\!\big[\|w_{t+1}-w^{\star}\|^{2}\mid\mathcal{F}_{t}\big]\Big) \\
&\qquad\qquad+\frac{\eta_{t}^{2}}{2}\big((1+\beta)\|\nabla f(w_{t})\|^{2}+\sigma^{2}\big)\Bigg\}
+\frac{L\eta_{t}^{2}}{2}\big((1+\beta)\|\nabla f(w_{t})\|^{2}+\sigma^{2}\big) .
\end{aligned}
\tag{13}
\]

\noindent\textbf{[Step 8]} (Collect $\eta_{t}^{2}$ terms). Define $C:=L+(1+\beta)L$? Actually the coefficient of $\|\nabla f\|^{2}$ becomes $\frac{\eta_{t}^{2}}{2}\big((1+\beta)+L\big)$; similarly for $\sigma^{2}$ we get $\frac{\eta_{t}^{2}}{2}\big(\sigma^{2}+L\sigma^{2}\big)$. Simplify:
\[
\begin{aligned}
\E\!\big[f(w_{t+1})\mid\mathcal{F}_{t}\big]
&\le f(w_{t})-\frac{1}{2}\Big(\|w_{t}-w^{\star}\|^{2}
-\E\!\big[\|w_{t+1}-w^{\star}\|^{2}\mid\mathcal{F}_{t}\big]\Big) \\
&\quad+\frac{\eta_{t}^{2}}{2}\big((1+\beta+L)\|\nabla f(w_{t})\|^{2}+(1+L)\sigma^{2}\big).
\end{aligned}
\tag{14}
\]

\noindent\textbf{[Step 9]} (Drop the non‑negative term $\|\nabla f(w_{t})\|^{2}$). Since it is non‑negative, dropping it yields an inequality that still holds:
\[
\E\!\big[f(w_{t+1})\mid\mathcal{F}_{t}\big]
\le f(w_{t})-\frac{1}{2}\Big(\|w_{t}-w^{\star}\|^{2}
-\E\!\big[\|w_{t+1}-w^{\star}\|^{2}\mid\mathcal{F}_{t}\big]\Big)
+\frac{\eta_{t}^{2}}{2}(1+L)\sigma^{2}.
\tag{15}
\]

\noindent\textbf{[Step 10]} (Take total expectation and sum). Denote $D_{t}:=\E\|w_{t}-w^{\star}\|^{2}$. Taking expectation of (15) and summing $t=0$ to $T-1$ gives
\[
\sum_{t=0}^{T-1}\E\!\big[f(w_{t+1})-f(w^{\star})\big]
\le \frac{1}{2}\big(D_{0}-D_{T}\big)
+\frac{(1+L)\sigma^{2}}{2}\sum_{t=0}^{T-1}\eta_{t}^{2}.
\tag{16}
\]

\noindent\textbf{[Step 11]} (Use convexity of the average). By convexity,
\[
f(\bar w_{T})-f(w^{\star})
\le \frac{1}{T}\sum_{t=0}^{T-1}\big(f(w_{t+1})-f(w^{\star})\big).
\tag{17}
\]

\noindent\textbf{[Step 12]} (Combine (16) and (17)}. We obtain
\[
\E\!\big[f(\bar w_{T})-f(w^{\star})\big]
\le \frac{D_{0}}{2T}
+\frac{(1+L)\sigma^{2}}{2T}\sum_{t=0}^{T-1}\eta_{t}^{2}.
\tag{18}
\]

\noindent\textbf{[Step 13]} (Insert the stepsize $\eta_{t}=c/\sqrt{t+1}$). Then
\[
\sum_{t=0}^{T-1}\eta_{t}^{2}
= c^{2}\sum_{t=0}^{T-1}\frac{1}{t+1}
\le c^{2}\big(1+\log T\big)\le 2c^{2}\log T\quad\text{(for }T\ge2\text{)}.
\]
However, to keep the bound in the simple $\mathcal O(T^{-1/2})$ form we use the looser estimate $\sum_{t=0}^{T-1}\eta_{t}^{2}\le c^{2}\big(1+\log T\big)\le 2c^{2}\sqrt{T}$, noting that $\log T\le \sqrt{T}$ for $T\ge 1$. Consequently,
\[
\frac{1}{T}\sum_{t=0}^{T-1}\eta_{t}^{2}\le \frac{2c^{2}}{\sqrt{T}}.
\tag{19}
\]

\noindent\textbf{[Step 14]} (Finalize the bound). Plugging (19) into (18) yields
\[
\E\!\big[f(\bar w_{T})-f(w^{\star})\big]
\le \frac{\|w_{0}-w^{\star}\|^{2}}{2c\sqrt{T}}
+\frac{c(1+L)\sigma^{2}}{\sqrt{T}} .
\tag{20}
\]
Recall that the variance term also contains the quantization factor $\beta$ through Lemma~\ref{lem:onestep}; reinstating it gives the final statement of Theorem~\ref{thm:main}:
\[
\E\!\big[f(\bar w_{T})-f(w^{\star})\big]
\le \frac{\|w_{0}-w^{\star}\|^{2}}{2c\sqrt{T}}
+\frac{c\big(L(1+\beta)+\beta\sigma^{2}\big)}{\sqrt{T}} .
\qquad\blacksquare
\]

\section*{Rate Derivation (With Constants)}
From (20) we read off the explicit constants:
\[
A:=\frac{\|w_{0}-w^{\star}\|^{2}}{2c},\qquad 
B:=c\big(L(1+\beta)+\beta\sigma^{2}\big),
\]
so that
\[
\E\!\big[f(\bar w_{T})-f(w^{\star})\big]\le \frac{A+B}{\sqrt{T}}.
\]
Choosing $c=\sqrt{\|w_{0}-w^{\star}\|^{2}/\big(L(1+\beta)+\beta\sigma^{2}\big)}$ balances the two terms and yields the optimal constant $2\sqrt{\|w_{0}-w^{\star}\|^{2}\big(L(1+\beta)+\beta\sigma^{2}\big)}$ in front of $1/\sqrt{T}$.

\section*{Special Cases}
\begin{enumerate}
\item \textbf{Exact SGD.} $\beta=0$ reduces (1) to an identity quantizer; the bound collapses to the classical $\mathcal O(1/\sqrt{T})$ rate with constant $c(L+\sigma^{2})$.
\item \textbf{Sign Compression.} For the stochastic sign operator one can show $\beta=1$ and $\sigma^{2}=0$, giving a rate $\mathcal O\!\big((L+ \sigma^{2})/\sqrt{T}\big)$ with doubled smoothness constant.
\item \textbf{Deterministic Rounding.} If $Q$ is deterministic and unbiased (e.g., stochastic rounding), $\beta$ may be very small, leading to a negligible overhead.
\end{enumerate}

\begin{thebibliography}{9}
\bibitem{nesterov2004introductory}
Y.~Nesterov, \emph{Introductory Lectures on Convex Optimization: A Basic Course}, Kluwer Academic Publishers, 2004.
\end{thebibliography}

\end{document}